\documentclass[12pt]{article}

\usepackage{amssymb,amsmath,amsthm}

% gaya teorema (semuanya memakai pencacah yang sama)
\theoremstyle{plain}
\newtheorem{theorem}{Teorema}[section]
\newtheorem{lemma}[theorem]{Lema}
\newtheorem{proposition}[theorem]{Proposisi}
\newtheorem{corollary}[theorem]{Korolari}
\newtheorem{conjecture}[theorem]{Konjektur}
\newtheorem{exercise}[theorem]{Latihan}
\theoremstyle{definition}
\newtheorem{definition}[theorem]{Definisi}
\newtheorem{example}[theorem]{Contoh}
\theoremstyle{remark}
\newtheorem{remark}[theorem]{Catatan}

% beberapa pengaturan penomoran
\numberwithin{equation}{section}

\usepackage{fancyhdr}
\usepackage{lastpage}
\setlength{\headheight}{15.2pt}
\renewcommand{\footrulewidth}{0.4pt}% nilai bakunya 0pt
\setlength{\footskip}{30pt}
\pagestyle{fancy}

\makeatletter
\let\ps@plain\ps@fancy 
\makeatother

\lhead{MATH 742}
\chead{Aljabar Multilinear}
\rhead{Musim Semi 2023}
\lfoot{Terakhir Direvisi: \today}
\cfoot{}
\rfoot{\thepage\ dari \pageref{LastPage} }

\begin{document}

%%%%%%%%%%%%%%%%%%
%%%%%%%%%%%%%%%%%%
\title{Aljabar Multilinear}
\author{Alexander Duncan}

\maketitle

Di sini kita meninjau dan/atau memperkenalkan fakta-fakta baku dari aljabar
(multi)linear. Sebagian besar materi ini semestinya sudah dijumpai dalam mata
kuliah aljabar linear tingkat sarjana atau rangkaian ujian kualifikasi aljabar
baku, tetapi beberapa gagasan (seperti hasil kali tensor) mungkin belum.

Dua rujukan baku untuk aljabar tingkat pascasarjana ialah
\cite{DF} dan \cite{Lang}. Kita berfokus pada kasus ketika gelanggang dasar
merupakan medan, karena inilah penerapan utama kita.
Banyak gagasan ini dapat digeneralisasi atau diperluas ke gelanggang sebarang,
tetapi kita tidak melakukannya di sini.

\section{Endomorfisme}

Misalkan $k$ suatu medan.

\begin{definition}
Sebuah \emph{$k$-aljabar} $A$ sekaligus merupakan ruang vektor atas $k$ dan
gelanggang dengan satuan sedemikian sehingga $r(ab)=(ra)b=a(rb)$ untuk semua $r \in k$ dan $a,b \in A$.
Sebuah \emph{morfisme $k$-aljabar} $f \colon A \to B$ sekaligus merupakan
transformasi linear atas $k$ dan homomorfisme gelanggang.
\end{definition}

Definisi alternatif (yang ekuivalen) bagi $k$-aljabar ialah gelanggang
dengan satuan $A$ beserta homomorfisme gelanggang $\pi : k \to A$,
yang disebut \emph{morfisme struktur}, dengan citra termuat dalam
pusat $A$. Dengan definisi ini,
morfisme $k$-aljabar adalah homomorfisme gelanggang $f\colon A \to B$ sedemikian
sehingga $\pi_B = f \circ \pi_A$, dengan $\pi_A$ dan $\pi_B$ masing-masing
merupakan morfisme struktur dari $A$ dan $B$.

\begin{exercise}
Buktikan bahwa kedua definisi $k$-aljabar (berturut-turut, morfisme $k$-aljabar)
ekuivalen.
\end{exercise}

\begin{definition}
Misalkan $\operatorname{M}_{n,m}(k)$ adalah himpunan matriks
$n \times m$ dengan koefisien di $k$, dan
misalkan $\operatorname{M}_{n}(k) = \operatorname{M}_{n,n}(k)$. 
\end{definition}

Himpunan matriks $\operatorname{M}_{n,m}(k)$ merupakan ruang vektor atas $k$,
dan himpunan matriks persegi $\operatorname{M}_n(k)$ merupakan $k$-aljabar
terhadap perkalian matriks.

\begin{definition}
Untuk ruang vektor atas $k$, yaitu $V,W$, misalkan
$\operatorname{Hom}_k(V,W)$ adalah himpunan transformasi linear atas $k$
dari $V$ ke $W$.
Kita tulis $\operatorname{End}_k(V) = \operatorname{Hom}_k(V,V)$
untuk himpunan \emph{endomorfisme} dari $V$.
Kita tulis $V^\vee = \operatorname{Hom}_k(V,k)$
untuk \emph{ruang dual} dari $V$.
\end{definition}

Himpunan $\operatorname{Hom}_k(V,W)$ merupakan ruang vektor atas $k$, sedangkan
himpunan $\operatorname{End}_k(V)$ juga merupakan $k$-aljabar terhadap komposisi.

Diberikan matriks $A \in \operatorname{M}_{n,m}(k)$, kita memperoleh
dua peta kanonik.
Yaitu, kita mempunyai \emph{perkalian kiri}
$L_A \in \operatorname{Hom}(k^m,k^n)$ melalui $L_A(v)=Av$
dengan memandang $v$ sebagai vektor kolom,
dan \emph{perkalian kanan}
$R_A \in \operatorname{Hom}(k^n,k^m)$ melalui $R_A(v)=vA$
dengan memandang $v$ sebagai vektor baris.

\begin{theorem}
Misalkan $\psi: k^n \to V$ dan $\phi : k^m \to W$ merupakan isomorfisme
ruang vektor.
Maka terdapat isomorfisme ruang vektor
$\operatorname{M}_{n,m}(k) \to \operatorname{Hom}_k(V,W)$ yang diberikan oleh
$A \mapsto \phi \circ L_A \circ \psi^{-1}$.
Dalam kasus $W=V$ dan $\phi=\psi$, ini menjadi isomorfisme $k$-aljabar
$\operatorname{M}_n(k) \to \operatorname{End}_k(V)$.
\end{theorem}

\subsection{Ruang Dual}

Jika $V$ adalah ruang vektor dengan basis hingga
$\beta_1, \ldots, \beta_n$, maka terdapat \emph{basis dual} tunggal
$\beta^\vee_1, \ldots, \beta^\vee_n$
untuk $V^\vee$ sedemikian sehingga $\beta_i^\vee(\beta_j)=\delta_{ij}$,
dengan $\delta_{ij}$ menyatakan delta Kronecker.

Khususnya, untuk ruang vektor berdimensi hingga $V$,
terdapat isomorfisme \emph{nonkanonik} $V \cong V^\vee$,
yang bergantung pada pilihan basis.
Sebaliknya, terdapat isomorfisme \emph{kanonik}
$V \cong (V^\vee)^\vee$ dengan memetakan vektor $v \in V$ ke
fungsional “evaluasi” $\operatorname{ev}_v : f \mapsto f(v)$. 

\subsection{Polinomial Karakteristik}

\begin{definition}
Diberikan matriks persegi $A \in \operatorname{M}_n(k)$, kita definisikan
\emph{polinomial karakteristik}:
\[
\chi_A(t) = \det(tI-A)
\]
yang merupakan unsur $k[t]$.
Kita memakai konvensi tanda yang memastikan bahwa
polinomial karakteristik selalu monik.

\emph{Polinomial minimal} $m_A(t) \in k[t]$ adalah pembangkit monik
dari ideal utama $\ker(\psi_A)$, dengan
$\psi_A : k[t] \to \operatorname{M}_n(k)$ merupakan homomorfisme $k$-aljabar
yang ditentukan oleh $\psi_A(t)=A$.
\end{definition}

\begin{theorem}[Cayley-Hamilton]
$m_A(t)$ membagi $\chi_A(t)$
\end{theorem}

\begin{definition}
Multihimpunan \emph{nilai eigen} dari $A \in \operatorname{M}_n(k)$ adalah
multihimpunan akar $\chi_A(t)$ atas penutupan aljabar
$\overline{k}$ dari $k$.
\end{definition}

Jika $\lambda_1, \ldots, \lambda_n$ adalah nilai-nilai eigen dari $\chi_A(t)$,
dengan memperhitungkan multiplisitasnya,
maka kita mempunyai
\[
\chi_A(t) = t^n - e_1 t^{n-1} + e_2 t^{n-2} - e_3 t^{n-3} + \cdots \pm e_n
\]
dengan $e_1,\ldots, e_n$ merupakan \emph{fungsi simetris elementer}
dalam $\lambda_1,\ldots, \lambda_n$.

Sebagai rumus, fungsi simetris elementer $e_i$
diberikan oleh
\[
e_i = \sum_{X \in S(n,k)} \left(\prod_{i \in X} \lambda_i\right)
\]
dengan $S(n,k)$ merupakan himpunan semua subhimpunan $\{1,\ldots,n\}$ yang
berukuran $k$.

Kasus-kasus khusus yang penting ialah teras
\[
\operatorname{tr}(A) = e_1 = \lambda_1 + \ldots + \lambda_n
\]
dan determinan
\[
\det(A) = e_n = \lambda_1 \lambda_2 \cdots \lambda_{n-1} \lambda_n .
\]

Ingat bahwa dua matriks $A,B \in \operatorname{M}_n(k)$ disebut \emph{serupa}
jika terdapat matriks $P$ sedemikian sehingga $A=PBP^{-1}$.
Diberikan endomorfisme $f$ dari ruang vektor berdimensi $n$, yaitu $V$,
kita memperoleh matriks-matriks representasi yang berbeda bagi $f$, bergantung
pada pilihan basis. Namun, himpunan matriks yang merepresentasikan $f$
merupakan satu kelas keserupaan dalam $\operatorname{M}_n(k)$.

Polinomial minimal, polinomial karakteristik, nilai-nilai eigen,
teras, dan determinan semuanya invarian terhadap keserupaan.
Dengan demikian, kita dapat mendefinisikan polinomial minimal,
polinomial karakteristik, nilai-nilai eigen, teras, dan determinan suatu
endomorfisme secara kanonik tanpa bergantung pada pilihan
basis.

\subsection{Bentuk Kanonik Jordan}

\emph{Blok Jordan berukuran $n$ dengan nilai eigen $\lambda$} adalah
matriks berdimensi $n$
\[
J_n(\lambda) = \begin{pmatrix}
\lambda & 1 & 0 & \cdots & 0 & 0\\
0 & \lambda & 1 & \cdots & 0 & 0\\
0 & 0& \lambda & \cdots & 0 & 0\\
\vdots & \vdots & \vdots & \ddots & \vdots & \vdots \\
0 & 0 & 0 & \cdots & \lambda & 1\\
0 & 0 & 0 & \cdots & 0 & \lambda\\
 \end{pmatrix} .
\]

\begin{definition}
Matriks $A$ berada dalam \emph{Bentuk Kanonik Jordan} jika $A$ berbentuk
diagonal blok dengan blok-bloknya berupa blok Jordan:
\[
A = \begin{pmatrix}
J_{n_1}(\lambda_1) & \cdots & 0\\
\vdots & \ddots & \vdots\\
0 & \cdots & J_{n_r}(\lambda_r)\\
\end{pmatrix} 
\]
dengan $n_1,\ldots,n_r$ dan $\lambda_1,\ldots,\lambda_r$
tidak harus berbeda.
\end{definition}

Perhatikan bahwa \emph{matriks diagonal} tepat merupakan matriks dalam bentuk
kanonik Jordan dengan $n_1=\ldots=n_r=1$.
Suatu matriks disebut \emph{dapat didiagonalkan} jika serupa dengan matriks diagonal.

\begin{theorem}
Misalkan $A$ adalah matriks persegi yang nilai-nilai eigennya terdefinisi atas $k$.
Maka $A$ serupa dengan suatu matriks dalam Bentuk Kanonik Jordan,
yang tunggal hingga pengurutan ulang blok-blok Jordan.
\end{theorem}

Perhatikan bahwa syarat dalam teorema di atas selalu berlaku ketika
medan $k$ tertutup secara aljabar.

\begin{exercise}
\[J_n(\lambda)^m =
\begin{pmatrix}
\lambda^m & \binom{m}{1}\lambda^{m-1} & \binom{m}{2}\lambda^{m-2}
& \cdots & \binom{m}{n-2}\lambda^{m-n-2} & \binom{m}{n-1}\lambda^{m-n-1}\\
0 & \lambda^m & \binom{m}{1}\lambda^{m-1}
& \cdots & \binom{m}{n-3}\lambda^{m-n-3} & \binom{m}{n-2}\lambda^{m-n-2}\\
0 & 0 & \lambda^m & \cdots &
\binom{m}{n-4}\lambda^{m-n-4} & \binom{m}{n-3}\lambda^{m-n-3}\\
\vdots & \vdots & \vdots & \ddots & \vdots & \vdots \\
0 & 0 & 0 & \cdots & \lambda^m & \binom{m}{1}\lambda^{m-1}\\
0 & 0 & 0 & \cdots & 0 & \lambda^m\\
 \end{pmatrix} .\]
\end{exercise}

\begin{corollary}
Misalkan $k$ adalah medan tertutup secara aljabar,
dan $m$ merupakan bilangan bulat yang relatif prima terhadap karakteristiknya.
Jika $A \in \operatorname{M}_n(k)$ berorde $m$,
maka $A$ dapat didiagonalkan dan nilai-nilai eigennya merupakan akar kesatuan ke-$m$.
\end{corollary}

\begin{proof}
Orde suatu matriks merupakan invarian keserupaan, karena
$A=PBP^{-1}$ mengakibatkan $A^m=PB^mP^{-1}$.
Jadi, kita dapat mengasumsikan bahwa $A$ berada dalam Bentuk Kanonik Jordan.
Karena kita mengetahui $A^m=I$, dari latihan tersebut kita menyimpulkan bahwa
untuk setiap nilai eigen $\lambda_i$ dari $A$, berlaku
$\lambda_i^m = 1$,
dan selain itu $m\lambda_i^{m-1}=0$ jika blok Jordan yang bersesuaian
bukan berukuran $1\times 1$.
Syarat pertama memaksa semua nilai eigen menjadi akar kesatuan ke-$m$,
sedangkan syarat kedua memaksa semua blok berukuran $1 \times 1$
(karena $m\ne 0$).
\end{proof}

Syarat bahwa bilangan bulat $m$ relatif prima terhadap karakteristik
memang diperlukan. Jika $k$ berkarakteristik $p$, maka matriks
\[
\begin{pmatrix} 1 & 1 \\ 0 & 1 \end{pmatrix} 
\]
berorde $p$, tetapi tidak dapat didiagonalkan. 

\subsection{Bentuk Kanonik Rasional}

TODO

\section{Peta Bilinear}

\begin{definition}
Misalkan $U$, $V$, dan $W$ adalah ruang vektor atas $k$.
Sebuah \emph{peta bilinear} $b: V \times W \to U$
adalah fungsi sedemikian sehingga
\begin{itemize}
\item $b(\lambda v_1+v_2,w) = \lambda b(v_1,w) + b(v_2,w)$, dan
\item $b(v,\lambda w_1+w_2) = \lambda b(v,w_1) + b(v,w_2)$
\end{itemize}
untuk semua $v,v_1,v_2 \in V$, $w,w_1,w_2 \in W$, dan $\lambda \in k$.
\end{definition}

Himpunan peta bilinear $\operatorname{Bil}_k(V,W;U)$ merupakan ruang vektor atas $k$,
yaitu ruang yang terdiri atas peta-peta bilinear tersebut.
Kita memakai singkatan $\operatorname{Bil}_k(V,W)=\operatorname{Bil}_k(V,W;k)$
jika kodomainnya adalah $k$.
Kita mempunyai isomorfisme kanonik
\[
\operatorname{Bil}_k(V,W) \cong
\operatorname{Hom}_k(V,W^\vee)
\]
dengan $b \in \operatorname{Bil}_k(V,W)$ dipetakan ke $b_1 : V \to W^\vee$
yang didefinisikan oleh $b_1(v) = b(v,-)$ untuk setiap $v \in V$.
Serupa dengan itu, kita mempunyai isomorfisme kanonik
\[
\operatorname{Bil}_k(V,W) \cong
\operatorname{Hom}_k(W,V^\vee)
\]
dengan $b \in \operatorname{Bil}_k(V,W)$ dipetakan ke $b_2 : W \to V^\vee$
yang didefinisikan oleh $b_2(w) = b(-,w)$ untuk setiap $w \in W$.

Misalkan $V$ mempunyai basis $v_1,\ldots,v_n$ dan $W$ mempunyai basis $w_1,\ldots, w_m$,
maka kita dapat membentuk \emph{matriks} $B$ dari peta bilinear $b$ dengan mendefinisikan
$B_{ij} = b(v_i,w_i)$.

\begin{exercise}
Pernyataan-pernyataan berikut ekuivalen:
\begin{itemize}
\item $b_1 : W \to V^\vee$ merupakan isomorfisme.
\item $b_2 : V \to W^\vee$ merupakan isomorfisme.
\item $B$ merupakan matriks invertibel (apa pun basis yang dipilih).
\end{itemize}
Jika salah satu syarat ini berlaku, kita katakan bahwa peta bilinear tersebut
\emph{tak degenerat}.
\end{exercise}

Peta bilinear tak degenerat $V \times W \to k$ sering disebut
\emph{pasangan sempurna}. Berdasarkan latihan tersebut, pasangan sempurna
menginduksi isomorfisme $V \cong W^\vee$ dan $W \cong V^\vee$.

\subsection{Bentuk Bilinear}

Sebuah \emph{bentuk bilinear} pada $V$ adalah peta bilinear $b : V \times V \to k$.
Dengan kata lain, kedua komponen hasil kali dalam domain merupakan
ruang vektor yang sama.
Dalam kasus ini, \emph{matriks} $B$ dari bentuk bilinear $b$ biasanya dibentuk
terhadap basis yang sama pada kedua komponen hasil kali.

\begin{definition}
Sebuah bentuk bilinear $b: V \times V \to k$ disebut
\begin{itemize}
\item \emph{simetris} jika $b(v,w)=b(w,v)$ untuk semua $v,w \in V$.
\item \emph{simetris-miring} atau \emph{antisimetris}
jika $b(v,w)=-b(w,v)$ untuk semua $v,w \in V$.
\item \emph{berselang-seling} jika $b(v,v)=0$ untuk semua $v \in V$.
\end{itemize}
\end{definition}

\begin{exercise}
Himpunan bentuk bilinear simetris (berturut-turut simetris-miring dan
berselang-seling) merupakan subruang dari $\operatorname{Bil}_k(V,V)$.
\end{exercise}

Jika $b: V \times V \to k$ merupakan bentuk bilinear simetris,
maka kedua peta kanonik $b_1 : V \to V^\vee$ dan $b_2 : V \to
V^\vee$ sama. Jadi, bentuk bilinear simetris menghasilkan
satu pilihan isomorfisme $V \to V^\vee$ (alih-alih sepasang).

\begin{example}
\emph{Hasil kali titik} pada $k^n$ merupakan bentuk bilinear simetris.
Dalam kasus $k$ tidak berkarakteristik $2$, bentuk tersebut
tak degenerat, dan isomorfisme yang bersesuaian
$k^n \to (k^n)^\vee$ dapat dipandang sebagai operasi
transpos, yang memetakan vektor kolom menjadi vektor baris dan
sebaliknya.
\end{example}

\begin{example}
Perhatikan matriks blok berukuran $2n \times 2n$
\[
J = \begin{pmatrix} 0 & I_n \\ -I_n & 0 \end{pmatrix} 
\]
dalam $k^{2n}$.
Bentuk bilinear yang bersesuaian merupakan bentuk bilinear
berselang-seling yang tak degenerat.
\end{example}

\begin{example}
Jika $k$ berkarakteristik $2$, maka “hasil kali titik”
pada $k^n$ merupakan bentuk bilinear simetris yang \emph{degenerat}.
Ini juga merupakan contoh bentuk bilinear simetris-miring yang tidak
berselang-seling.
\end{example}

\begin{proposition}
Setiap bentuk berselang-seling bersifat simetris-miring.
\end{proposition}

\begin{proof}
Jika $b$ berselang-seling, maka
\begin{align*}
0 &= b(v+w,v+w)\\
&= b(v,v) + b(v,w) + b(w,v) + b(w,w)\\
&= b(v,w)+b(w,v)
.\end{align*}
Jadi, kita menyimpulkan bahwa $b$ bersifat simetris-miring.
\end{proof}

\begin{proposition}
Jika $k$ \emph{tidak} berkarakteristik $2$, maka suatu bentuk bilinear
berselang-seling jika dan hanya jika bentuk itu simetris-miring.
\end{proposition}

\begin{proof}
Jika $b$ simetris-miring, maka $b(v,v)=-b(v,v)$ untuk semua $v \in V$.
Dengan demikian $2b(v,v)=0$ dan, karena $\frac{1}{2} \in k$,
kita melihat bahwa $b$ berselang-seling.
\end{proof}

\begin{proposition}
Jika $k$ berkarakteristik $2$, maka suatu bentuk bilinear
simetris jika dan hanya jika bentuk itu simetris-miring.
Khususnya, setiap bentuk berselang-seling bersifat simetris.
\end{proposition}

\begin{proof}
Jelas.
\end{proof}

\subsection{Bentuk Sesquilinear}

Dalam subbagian ini, kita bekerja atas bilangan kompleks.
Gagasan-gagasan ini dapat diperluas ke perluasan kuadratik dari medan-medan
lain, tetapi kita tidak akan membahasnya terlalu jauh.
Di bawah ini, jika $z=a+bi \in \mathbb{C}$, kita tulis
$\overline{z}=a-bi$ untuk konjugat kompleksnya.

\begin{definition}
Misalkan $U$, $V$, dan $W$ adalah ruang vektor atas $\mathbb{C}$.
Sebuah \emph{peta sesquilinear} $s: V \times W \to U$
adalah fungsi sedemikian sehingga
\begin{itemize}
\item $s(\lambda v_1+v_2,w) = \overline{\lambda} s(v_1,w) + s(v_2,w)$, dan
\item $s(v,\lambda w_1+w_2) = \lambda s(v,w_1) + s(v,w_2)$
\end{itemize}
untuk semua $v,v_1,v_2 \in V$, $w,w_1,w_2 \in W$, dan $\lambda \in \mathbb{C}$.
\end{definition}

Peta sesquilinear \emph{bukan} bentuk bilinear, karena peta itu linear
konjugat pada entri pertama, bukan linear.
Kita memakai “konvensi fisika”, yaitu entri pertama bersifat linear
konjugat, bukan entri kedua. Ini berbeda dari konvensi dalam banyak bidang
matematika, yang menetapkan bahwa entri \emph{kedua} bersifat linear konjugat.
Konvensi pertama jauh lebih alami menurut konvensi matriks baku
dan tampaknya makin banyak dipakai dalam teks-teks modern.

\begin{definition}
Sebuah \emph{bentuk sesquilinear} adalah peta sesquilinear $s: V \times V \to
\mathbb{C}$.
Bentuk sesquilinear disebut \emph{Hermitian} jika $s(v,w)=\overline{s(w,v)}$
untuk semua $v,w \in V$.
Bentuk Hermitian disebut \emph{definit positif} jika $s(v,v) > 0$ jika dan hanya
jika $v \ne 0$.
\end{definition}

Secara informal, jika kita membatasi bentuk sesquilinear pada ruang vektor atas
submedan real total $k \subseteq \mathbb{R}$, maka sifat Hermitian membatasi menjadi simetris,
dan sifat definit positif membatasi menjadi tak degenerat.

\subsection{Grup Klasik}

TODO

\subsection{Peta Multilinear}

\begin{definition}
Misalkan $U,V_1,\ldots, V_d$ adalah ruang vektor atas $k$.
Sebuah \emph{peta multilinear} $m: V_1 \times \cdots \times V_d \to U$
adalah fungsi yang linear pada setiap entri; dengan kata lain,
\begin{align*}
&m(v_1,\ldots,v_{i-1},\lambda v_i + v_i',v_{i+1},\ldots,v_d)\\
=& \lambda m(v_1,\ldots,v_{i-1},v_i,v_{i+1},\ldots,v_d)
+ m(v_1,\ldots,v_{i-1},v_i',v_{i+1},\ldots,v_d)
\end{align*}
untuk setiap indeks $i$ dan semua $v_j,v_j' \in V_j$, $\lambda \in k$.
Sebuah \emph{bentuk multilinear} adalah peta multilinear dengan $V_1 = \cdots = V_d$.
\end{definition}

\begin{definition}
Sebuah bentuk multilinear disebut
\emph{simetris} jika
\[
m(\cdots,v_i,\cdots,v_j,\cdots) = m(\cdots,v_j,\cdots,v_i,\cdots)
\]
untuk semua pasangan indeks $i,j$.
\end{definition}

\begin{definition}
Sebuah bentuk multilinear disebut \emph{simetris-miring} atau \emph{antisimetris} jika
\[
m(\cdots,v_i,\cdots,v_j,\cdots) = -m(\cdots,v_j,\cdots,v_i,\cdots)
\]
untuk semua pasangan indeks berbeda $i,j$.
\end{definition}

\begin{definition}
Sebuah bentuk multilinear disebut \emph{berselang-seling} jika
\[
m(v_1,\ldots,v_d) = 0
\]
setiap kali $v_i=v_j$ untuk suatu $i \ne j$.
\end{definition}

Seperti dalam kasus bilinear, peta-peta multilinear membentuk ruang vektor,
sedangkan bentuk-bentuk simetris, simetris-miring, dan berselang-seling
merupakan subruang dari ruang bentuk multilinear.

\begin{exercise}
Buktikan bahwa
\begin{itemize}
\item bentuk berselang-seling bersifat simetris-miring,
\item jika $\operatorname{char}(k) \ne 2$, maka bentuk simetris-miring
berselang-seling, dan
\item jika $\operatorname{char}(k) = 2$, maka bentuk simetris-miring
simetris, dan sebaliknya.
\end{itemize}
\end{exercise}

\begin{example}
Determinan $\det : M_n(k) \to k$ merupakan bentuk multilinear berselang-seling
pada baris-baris (atau kolom-kolom) matriks.
Selain itu, determinan merupakan satu-satunya bentuk demikian yang bernilai $1$ pada matriks identitas.
\end{example}

\section{Hasil Kali Tensor}

Pembaca dapat merujuk \cite[10.4]{DF} atau \cite[XVI]{Lang}
untuk uraian lebih lanjut tentang sifat-sifat hasil kali tensor dalam konteks
teori modul yang umum.
Di sini kita berfokus pada kasus ruang vektor berdimensi hingga,
yang sedikit lebih mudah.

Kita mendefinisikan hasil kali tensor melalui sifat pemetaan universalnya.

\begin{definition}
Misalkan $V_1,\ldots, V_n$ adalah ruang vektor atas $k$.
Sebuah \emph{hasil kali tensor} $V_1 \otimes_k V_2 \otimes \cdots \otimes_k V_d$
adalah ruang vektor atas $k$ beserta peta multilinear
\[
\pi : V_1 \times \cdots \times V_d
\to  V_1 \otimes_k \cdots \otimes_k V_d
\]
sedemikian sehingga untuk setiap peta multilinear $\phi : V_1 \times \cdots \times V_d
\to W$ terdapat peta linear tunggal
$\psi : V_1 \otimes_k \cdots \otimes_k V_d \to W$ sedemikian sehingga
$\phi = \psi \circ \pi$.
\end{definition}

Jika medan dasar $k$ sudah jelas, kita sering menghilangkannya dari
notasi; dengan kata lain, $V \otimes U$ merupakan singkatan bagi $V \otimes_k U$.
Kita juga dapat memakai notasi singkat
\[
\bigotimes_{i=1}^d V_i := V_1 \otimes \cdots \otimes V_d .
\]

Diberikan $(v_1,\ldots, v_d) \in V_1 \times \cdots \times V_d$,
misalkan $v_1 \otimes v_2 \otimes \cdots \otimes v_d$ menyatakan citra
$\pi(v_1,\ldots,v_d)$ dalam $V_1 \otimes_k \cdots \otimes_k V_d$.
Unsur-unsur ini kita sebut \emph{tensor sederhana}.

\begin{theorem}
Hasil kali tensor ada dan tunggal hingga isomorfisme tunggal.
\end{theorem}

\begin{proof}
Kita uraikan suatu konstruksi dan menyerahkan perincian sisanya sebagai latihan.
Misalkan $\Omega$ adalah ruang vektor bebas atas $k$ dengan basis berupa unsur-unsur
$V_1 \times \cdots \times V_d$. Misalkan $[x]$ adalah unsur basis
yang bersesuaian dengan $x$.
Misalkan $I$ adalah subruang $\Omega$ yang direntang oleh
\[
[(v_1, \ldots, (\lambda v_i + v'_i), \ldots, v_d)]
-
\lambda[(v_1, \ldots,v_i, \ldots, v_d)] -
[(v_1, \ldots, v'_i, \ldots, v_d)]
\]
untuk setiap indeks $i$, setiap $\lambda \in k$, dan semua $v_j,v'_j \in V_j$
untuk setiap $j$.
Definisikan $V_1 \otimes \cdots \otimes V_n$ sebagai ruang hasil bagi
$\Omega/I$, dengan $\pi$ berupa komposisi inklusi kanonik
$\prod_{i=1}^d V_i \to \Omega$ dan peta hasil bagi $\Omega \to
\bigotimes_{i=1}^d V_i$.
\end{proof}

Kita memperoleh korolari-korolari berikut:

\begin{corollary}
Ruang vektor $V_1 \otimes_k \cdots \otimes_k V_d$
direntang oleh tensor-tensor sederhananya.
Relasi
\begin{align*}
&v_1 \otimes \cdots \otimes (\lambda v_i + v'_i) \otimes \cdots \otimes
v_d\\
=&
\lambda(v_1 \otimes \cdots \otimes v_i \otimes \cdots \otimes v_d) +
v_1 \otimes \cdots \otimes v'_i \otimes \cdots \otimes v_d
\end{align*}
berlaku
untuk setiap indeks $i$, setiap $\lambda \in k$, dan semua $v_j,v'_j \in V_j$.
\end{corollary}

\begin{proof}
Latihan.
\end{proof}

\begin{proposition}
Jika $B_1,\ldots B_d$ masing-masing merupakan basis bagi $V_1,\ldots,V_d$, maka
\[
\{ b_1 \otimes \cdots \otimes b_d \mid b_1 \in B_1, \ldots, b_d \in B_d
\}
\]
merupakan basis bagi $V_1 \otimes_k \cdots \otimes_k V_d$.
\end{proposition}

\begin{proof}
Latihan.
\end{proof}

\begin{remark}
Transformasi linear sering didefinisikan dengan
“memperluas secara linear”.
Misalkan kita mempunyai ruang vektor $V,W$ dan himpunan perentang
$S$ bagi $V$. Jika kita mendefinisikan fungsi $f : S \to W$,
maka paling banyak terdapat satu transformasi linear $F : V \to W$
yang memperluas $f$. Jika $S$ merupakan \emph{basis}, perluasan $F$ selalu ada,
tetapi jika $S$ tidak bebas linear, kita harus memastikan bahwa setiap
perluasan linear terdefinisi dengan baik.
Dalam kasus hasil kali tensor, himpunan $S$ sering diambil sebagai
himpunan tensor sederhana.
\end{remark}

\begin{theorem}
Untuk ruang vektor berdimensi hingga $V$ dan ruang vektor sebarang
$W$, kita mempunyai isomorfisme kanonik
\[
V^\vee \otimes W \cong \operatorname{Hom}_k(V,W) .
\]
\end{theorem}

\begin{proof}
Definisikan $\psi : V^\vee \otimes W \cong \operatorname{Hom}_k(V,W)$
dengan menetapkan
\[
\psi(f \otimes w)(v) = f(v)w
\]
untuk $f \in V^\vee$, $w \in W$, $v \in V$, lalu memperluas secara linear.
(Latihan: periksa bahwa perluasan secara linear sah di sini!)

Kita mengonstruksi peta invers $\phi$ secara eksplisit.
Misalkan $v_1,\ldots, v_d$ adalah basis bagi $V$.
Diberikan peta linear $g : V \to W$, definisikan
\[
\phi(g) = \sum_{i=1}^d v^\vee_i \otimes g(v_i) \ .
\]
(Latihan: periksa bahwa $\phi \circ \psi$ dan $\psi \circ \phi$
merupakan identitas.)
\end{proof}

\begin{exercise}
Dengan mengidentifikasi $V^\vee \otimes V \cong \operatorname{Hom}_k(V,V)$,
peta linear $T: V^\vee \otimes V \to k$ yang didefinisikan oleh
\[
T\left( \sum_{i=1}^d f_i \otimes v_i\right)
= \sum_{i=1}^d f_i(v_i)
\]
dapat diidentifikasi dengan teras $\operatorname{tr} :
\operatorname{Hom}_k(V,V) \to k$.
\end{exercise}

\begin{corollary}
Untuk ruang vektor berdimensi hingga $V, W$,
kita mempunyai isomorfisme kanonik
\[
V^\vee \otimes W^\vee \cong \operatorname{Bil}_k(V,W) .
\]
\end{corollary}

\begin{proof}
Keduanya secara kanonik isomorfik dengan $\operatorname{Hom}_k(V,W^\vee)$.
\end{proof}

\begin{proposition}
Misalkan $U,V,W$ adalah ruang vektor berdimensi hingga.
Terdapat isomorfisme-isomorfisme kanonik berikut:
\begin{enumerate}
\item $V \otimes W \cong W \otimes V$
\item $U \otimes (V \oplus W) \cong (U \otimes V) \oplus (U \otimes W)$
\item $U \otimes V \otimes W \cong U \otimes (V \otimes W)
\cong (U \otimes V) \otimes W$
\item $V \otimes k \cong V$
\item $(V \otimes W)^\vee \cong V\vee \otimes W^\vee$
\item $\operatorname{Hom}(U,\operatorname{Hom}(V,W)
\cong \operatorname{Hom}(U \otimes V, W)$
\end{enumerate}
\end{proposition}

\begin{proof}
Semua ini kita serahkan sebagai latihan.
Salah satu petunjuk yang berguna ialah bahwa
$(V \otimes W)^\vee = \operatorname{Hom}(V \otimes W,k)$
secara kanonik isomorfik dengan $\operatorname{Bil}(V,W)$
berdasarkan Sifat Pemetaan Universal hasil kali tensor.
\end{proof}

\begin{proposition}
Misalkan $L/k$ merupakan perluasan medan.
Jika $V$ adalah ruang vektor atas $k$, maka $V \otimes_k L$ merupakan ruang vektor atas $L$.
Jika $A$ adalah $k$-aljabar, maka $A \otimes_k L$ merupakan $L$-aljabar.
\end{proposition}

\begin{proof}
Himpunan $V \otimes_k L$ sudah merupakan grup abelian.
Kita mendefinisikan perkalian skalar dengan memperluas secara linear
rumus $\lambda(v\otimes l)=v \otimes \lambda l$
untuk $\lambda, l \in L$ dan $v \in V$.
Untuk melihat bahwa $A$ merupakan $L$-aljabar, kita perlu mendefinisikan perkalian gelanggangnya.
Kita menetapkan $(a_1 \otimes l_1)(a_2 \otimes l_2)=a_1a_2 \otimes l_1l_2$.
\end{proof}

\section{Aljabar Tensor}

Untuk sebagian materi ini, kita merujuk ke \cite[\S{11.5}]{DF} atau \cite{Lang}.
Materi dalam \cite[\S{A.2}]{Eisenbud} menyajikan sudut pandang alternatif
yang sedikit lebih umum dan abstrak daripada yang kita perlukan.

\begin{definition}
Sebuah \emph{$k$-aljabar berderajat} adalah $k$-aljabar $A$ dengan dekomposisi
jumlah langsung
\[
A = \bigoplus_{n \ge 0} A_n
\]
sedemikian sehingga $A_n \cdot A_m \subseteq A_{n+m}$ untuk semua bilangan bulat $n,m$.
Unsur $f \in A_d$ disebut \emph{homogen berderajat $d$}.
\end{definition}

Contoh baku $k$-aljabar berderajat ialah
$k$-aljabar polinomial $k[x_1,\ldots,x_n]$ yang diberi derajat menurut derajat polinomial.

\begin{definition}
Misalkan $V$ adalah ruang vektor berdimensi hingga atas $k$. Untuk bilangan bulat positif
$n$, kita definisikan $\mathcal{T}^d(V)=\bigotimes_{i=1}^d V$ dan
$\mathcal{T}^0(V)=k$. Kita definisikan \emph{$k$-aljabar tensor pada $V$}
sebagai $k$-aljabar berderajat
\[
\mathcal{T}(V) := \bigoplus_{d \ge 0} \mathcal{T}^d(V)
\]
dengan perkalian
\[\mathcal{T}^d(V) \times \mathcal{T}^e(V)
\to \mathcal{T}^{d+e}(V)
\]
yang didefinisikan melalui
\[
(v_1 \otimes \cdots \otimes v_d) (w_1 \otimes \cdots \otimes w_e)
=v_1 \otimes \cdots \otimes v_d \otimes w_1 \otimes \cdots \otimes w_e
\]
dan diperluas ke seluruh $\mathcal{T}(V)$ secara linear. 
\end{definition}

\begin{example}
Jika $V=k^n$, maka $\mathcal{T}(V)$ dapat diidentifikasi dengan
gelanggang polinomial nonkomutatif $k\langle x_1,\ldots, x_n\rangle$.
Sebagai contoh, jika $e_1,\ldots, e_n$ adalah basis baku bagi $k^n$,
maka
\[
e_1 \otimes e_1 \otimes e_1 + 3e_1 \otimes e_2 - e_2\otimes e_1 + 4\cdot 1
\]
dalam $\mathcal{T}(k^n)$
bersesuaian dengan unsur tak homogen
\[
x_1^3 + 3x_1x_2 - x_2x_1 + 4
\]
dalam $k\langle x_1,\ldots, x_n\rangle$.
\end{example}

Pernyataan berikut langsung diperoleh dari uraian di atas,
tetapi berguna untuk membandingkannya dengan aljabar simetris dan eksterior
yang dibahas di bawah.

\begin{proposition}
Jika $e_1, \ldots, e_n$ adalah basis bagi $V$, maka
\[
\{
e_{i_1} \otimes e_{i_2} \otimes \cdots \otimes e_{i_d}
\mid
(i_1,\ldots,i_d) \in \{1,\ldots,n\}^d
\}
\]
merupakan basis bagi $\mathcal{T}^d(V)$.
Khususnya, $\dim_k\ \mathcal{T}^d(k^n)=n^d$.
\end{proposition}

\subsection{Aljabar Simetris}

\begin{definition}
Misalkan $V$ adalah ruang vektor berdimensi hingga atas $k$.
\emph{Aljabar simetris pada $V$} adalah $k$-aljabar berderajat
\[
\mathcal{S}(V) = \mathcal{T}(V)/I
\]
dengan $I$ merupakan ideal berderajat (dua sisi)
\[
I = ( v \otimes w - w \otimes v \mid w,v \in V ).
\]
Komponen berderajat $\mathcal{S}^d(V)$ disebut
\emph{pangkat simetris ke-$d$}.
\end{definition}

Pangkat simetris memenuhi sifat pemetaan universal yang dibangun dari
sifat bagi hasil kali tensor:

\begin{proposition}
Komposisi
\[
\pi : V^d \to \mathcal{T}^d(V) \to \mathcal{S}^d(V)
\]
merupakan peta multilinear simetris sedemikian sehingga untuk setiap
peta multilinear simetris
$\phi : V^d \to W$ terdapat peta linear tunggal
$\psi : \mathcal{S}^d(V) \to W$
sedemikian sehingga
$\phi = \psi \circ \pi$.
\end{proposition}

Jika $V$ adalah ruang vektor berdimensi $n$, maka
$\mathcal{S}(V) \cong k[x_1,\ldots,x_n]$ sebagai $k$-aljabar berderajat.
Dengan demikian, aljabar simetris pada dasarnya merupakan versi “bebas koordinat” dari
gelanggang polinomial.
Karena alasan ini, perkalian dalam aljabar simetris biasanya cukup ditulis
sebagai monomial dalam vektor-vektor penyusunnya.
Sebagai contoh, citra $e_1 \otimes e_2 \otimes e_1$ dapat
ditulis $e_1^2e_2$.

Kita mempunyai proposisi kombinatorial berikut.

\begin{proposition}
Jika $e_1, \ldots, e_n$ adalah basis bagi $V$, maka
\[
\{
e_{i_1}e_{i_2} \cdots e_{i_d}
\mid
1 \le i_1 \le \cdots \le i_d \le n
\}
\]
merupakan basis bagi $\mathcal{S}^d(V)$.
Khususnya, $\dim_k\ \mathcal{S}^d(k^n)=\binom{n+d-1}{d}$.
\end{proposition}

\subsection{Aljabar Eksterior}

\begin{definition}
Misalkan $V$ adalah ruang vektor berdimensi hingga atas $k$.
\emph{Aljabar eksterior pada $V$} adalah $k$-aljabar berderajat
\[
\Lambda(V) = \mathcal{T}(V)/J
\]
dengan $J$ merupakan ideal berderajat (dua sisi)
\[
J = ( v \otimes v \mid v \in V ).
\]
Komponen berderajat $\Lambda^d(V)$ disebut
\emph{pangkat eksterior ke-$d$}.
\end{definition}

Pangkat eksterior memenuhi sifat pemetaan universal yang dibangun dari
sifat bagi hasil kali tensor:

\begin{proposition}
Komposisi
\[
\pi : V^d \to \mathcal{T}^d(V) \to \Lambda^d(V)
\]
merupakan peta multilinear berselang-seling sedemikian sehingga untuk setiap
peta multilinear berselang-seling
$\phi : V^d \to W$ terdapat peta linear tunggal
$\psi : \Lambda^d(V) \to W$
sedemikian sehingga
$\phi = \psi \circ \pi$.
\end{proposition}

Citra tensor sederhana $v_1 \otimes v_2 \otimes \cdots \otimes v_d$
ditulis $v_1 \wedge v_2 \wedge \cdots \wedge v_d$.
Demikian pula, hasil kali dua unsur $f,g \in \Lambda(V)$
ditulis $f \wedge g$.

\begin{proposition}
Jika $v,w \in V$, maka $v \wedge w = -w \wedge v$.
\end{proposition}

\begin{proof}
Ini langsung diperoleh dari
\begin{align*}
0 &= (v+w) \wedge (v+w)\\
&=v \wedge v + v \wedge w + w \wedge v + w \wedge w\\
&=v \wedge w + w \wedge v
\end{align*}
\end{proof}

\begin{proposition}
Jika $v_1,\ldots,v_d \in V$ dan $\sigma \in S_d$,
maka
\[v_1 \wedge \cdots \wedge v_d =
\operatorname{sgn}(\sigma)
v_{\sigma(1)} \wedge \cdots \wedge v_{\sigma(d)}.
\]
\end{proposition}

\begin{proof}
Setiap permutasi dapat ditulis sebagai hasil kali $m$ transposisi
berdampingan, dan menurut definisi $\operatorname{sgn}(\sigma)=(-1)^m$.
Jadi, hasil tersebut diperoleh dari proposisi sebelumnya.
\end{proof}

\begin{proposition}
Jika $f \in \Lambda^d(V)$ dan $g \in \Lambda^e(V)$,
maka $f \wedge g = (-1)^{de} g \wedge f$.
\end{proposition}

\begin{proof}
Ini diperoleh dari proposisi sebelumnya dengan memperhatikan bahwa
permutasi
\[
(1,\ldots,d+e) \mapsto (d+1,\ldots,d+e,1,\ldots,d)
\]
bertanda $(-1)^{de}$.
\end{proof}

\begin{proposition}
Vektor-vektor $v_1,\ldots , v_d \in V$ bergantung linear
jika dan hanya jika $v_1 \wedge \cdots \wedge v_d=0$.
\end{proposition}

\begin{proof}
Misalkan $v_1,\ldots, v_d$ bergantung linear.
Tanpa mengurangi keumuman, kita dapat mengasumsikan
\[
v_d = \alpha_1 v_1 + \cdots + \alpha_{d-1} v_{d-1}
\]
dengan $\alpha_1, \ldots, \alpha_{d-1} \in k$.
Dengan demikian,
\[
v_1 \wedge \cdots \wedge v_d =
\sum_{i=1}^{d-1} \alpha_i(v_1 \wedge \cdots \wedge v_{d-1} \wedge v_i).
\]
Karena $v_i$ muncul dua kali dalam setiap suku, ruas kanan bernilai $0$.

Misalkan $v_1,\ldots, v_d$ bebas linear.
Kita dapat melengkapinya menjadi basis $v_1,\ldots,v_n$ bagi $V$.
Misalkan $\beta : V \to k^n$ adalah peta koordinat yang bersesuaian.
Misalkan $f : V^d \to k$ adalah fungsi
\[
f(w_1,\ldots,w_d) = \det
\begin{pmatrix} \beta(w_1) & \cdots & \beta(w_d)
& e_{d+1} & \cdots & e_n
\end{pmatrix}
\]
dengan $\beta(w_i)$ dan $e_j$ dipandang sebagai kolom-kolom sebuah matriks.
Karena determinan merupakan peta multilinear berselang-seling pada kolom-kolom
matriks persegi, kita melihat bahwa $f$ merupakan peta multilinear berselang-seling
sedemikian sehingga $f(v_1,\ldots,v_d)=1$.
Berdasarkan sifat pemetaan universal, jika terdapat peta multilinear
berselang-seling sedemikian sehingga $f(v_1,\ldots,v_d) \ne 1$,
maka $v_1 \wedge \cdots \wedge v_d$ harus tak nol.
\end{proof}

Kita mempunyai proposisi kombinatorial berikut.

\begin{proposition}
Jika $e_1, \ldots, e_n$ adalah basis bagi $V$, maka
\[
\{
e_{i_1} \wedge \cdots e_{i_2} \cdots \wedge e_{i_d}
\mid
1 \le i_1 < \cdots < i_d \le n
\}
\]
merupakan basis bagi $\Lambda^d(V)$.
Khususnya, $\dim_k\ \Lambda^d(k^n)=\binom{n}{d}$.
\end{proposition}

\subsection{Subruang dari Pangkat Tensor}

Terdapat aksi kiri alami dari grup simetris $S_d$ pada
pangkat tensor $\mathcal{T}^d(V)$ yang diberikan oleh
\[
\sigma(v_1 \otimes v_2 \otimes \cdots \otimes v_n)
= v_{\sigma^{-1}(1)} \otimes v_{\sigma^{-1}(2)} \otimes \cdots
\otimes v_{\sigma^{-1}(n)}
\]
untuk $\sigma \in S_d$.
Inversnya penting untuk menjamin bahwa ini merupakan aksi \emph{kiri},
tetapi hal tersebut tidak terlalu relevan bagi pembahasan berikut.

\begin{definition}
Misalkan $V$ adalah ruang vektor atas $k$.
\emph{Ruang tensor-$d$ simetris pada $V$} adalah
\[
\operatorname{Sym}^d(V) = \operatorname{span}_k
\left\{ t \in \mathcal{T}^d(V) \mid \sigma(t)=t
\textrm{ untuk semua } t \in S_d  \right\} .
\]
\emph{Ruang tensor-$d$ berselang-seling pada $V$} adalah
\[
\operatorname{Alt}^d(V) = \operatorname{span}_k
\left\{ t \in \mathcal{T}^d(V) \mid
\sigma(t)=\operatorname{sgn}(\sigma)t
\textrm{ untuk semua } t \in S_d  \right\} .
\]
\end{definition}

Perhatikan bahwa meskipun $\mathcal{S}^d(V)$ dan $\Lambda^d(V)$ merupakan
\emph{hasil bagi} dari $\mathcal{T}^d(V)$,
ruang $\operatorname{Sym}^d(V)$
dan $\operatorname{Alt}^d(V)$
merupakan \emph{subruang}.

\begin{example}
Kita mempunyai
\[
e_1 \otimes e_2 \otimes e_2 + e_2 \otimes e_1 \otimes e_2
+ e_2 \otimes e_2 \otimes e_1 \in \operatorname{Sym}^d(k^5)
\]
dan
\[
3\left(e_1 \otimes e_2 - e_2 \otimes e_1\right)
+ 4\left( e_3 \otimes e_5 - e_5 \otimes e_3 \right) 
\in \operatorname{Alt}^d(k^7).
\]
Perhatikan bahwa grup simetris \emph{tidak} beraksi pada
ruang vektor $V$.
\end{example}

\begin{proposition} \label{prop:sym_alt_good_char}
Misalkan $\operatorname{char}(k) =0$ atau $\operatorname{char}(k) > d$.
Maka komposisi-komposisi
\begin{align*}
\operatorname{Sym}^d(V) &\to \mathcal{T}^d(V) \to \mathcal{S}^d(V)\\
\operatorname{Alt}^d(V) &\to \mathcal{T}^d(V) \to \Lambda^d(V)
\end{align*}
merupakan isomorfisme.
\end{proposition}

\begin{proof}
Kita mendefinisikan peta \emph{polarisasi}
$P : \mathcal{S}^d(V) \to \operatorname{Sym}^d(V)$ sebagai berikut.
Pertama-tama, perhatikan peta $p : \mathcal{T}^d \to \mathcal{T}^d$
yang didefinisikan melalui
\[
p(v_1 \otimes v_2 \otimes \cdots \otimes v_d)
= \frac{1}{d!}\sum_{\sigma \in S_d}
v_{\sigma(1)} \otimes v_{\sigma(1)} \otimes \cdots \otimes v_{\sigma(d)}
.
\]
Perhatikan bahwa $p$ merupakan proyeksi linear ke $\operatorname{Sym}^d(V)$
dengan kernel tepat $I \cap \mathcal{T}^d(V)$ dari definisi
$\mathcal{S}^d(V)$. Dengan demikian, kita memperoleh peta $P$ yang diinginkan.
Dapat diperiksa bahwa $P$ merupakan invers dari komposisi dalam pernyataan
proposisi.

Terdapat konstruksi serupa untuk $\operatorname{Alt}^d(V)$
dan $\Lambda^d(V)$.
\end{proof}

Isomorfisme-isomorfisme di atas digunakan secara luas, tetapi dapat
sangat menyesatkan. Sebagai contoh, isomorfisme-isomorfisme tersebut memberikan
struktur aljabar berderajat pada jumlah langsung tak hingga
$\bigoplus_{d\ge 0}\operatorname{Sym}^d(V)$,
tetapi perkaliannya \emph{tidak} sama dengan perkalian yang diwarisi dari
$\mathcal{T}$.

\begin{proposition}
$\operatorname{Sym}^d(V^\vee)$ dan $\left(\mathcal{S}^d(V)\right)^\vee$
secara kanonik isomorfik.
\end{proposition}

\begin{proof}
Berdasarkan sifat pemetaan universal hasil kali tensor,
$\mathcal{T}^d(V^\vee)$ secara kanonik isomorfik dengan ruang
peta $d$-linear $V^d \to k$.
Jadi, $\operatorname{Sym}^d(V^\vee)$ isomorfik dengan
ruang bentuk $d$-linear \emph{simetris} $V^d \to k$.
Bentuk-bentuk ini berkorespondensi bijektif secara kanonik dengan peta linear
$\mathcal{S}^d(V) \to k$ berdasarkan sifat universal
$\mathcal{S}^d(V)$. Dengan demikian, diperoleh isomorfisme tersebut.
\end{proof}

Tidak mungkin terdapat isomorfisme kanonik
$\operatorname{Alt}^d(V^\vee) \cong \left(\Lambda^d(V)\right)^\vee$
dalam semua kasus. Memang,
$\operatorname{Alt}^d(V^\vee) = \operatorname{Sym}^d(V^\vee)$
dalam karakteristik $2$, yang mungkin memiliki dimensi berbeda dari
$\Lambda^d(V)^\vee$.
Namun, pernyataan berikut berlaku dalam semua karakteristik.

\begin{proposition}
$\Lambda^d(V^\vee)$ dan $\left(\Lambda^d(V)\right)^\vee$
secara kanonik isomorfik.
\end{proposition}

\begin{proof}
Isomorfisme tersebut setara dengan pasangan sempurna
\[
\Lambda^d(V^\vee) \times \Lambda^d(V) \to k
\]
yang didefinisikan oleh
\[
(f_1 \wedge \cdots \wedge f_d)(v_1 \wedge \cdots \wedge v_d)
= \det \begin{pmatrix}
f_1(v_1) & \cdots & f_d(v_1)\\
\vdots & \ddots & \vdots\\
f_1(v_d) & \cdots & f_d(v_d)
\end{pmatrix} 
\]
untuk semua $v_1,\ldots, v_d \in V$ dan $f_1,\ldots, f_d \in V^\vee$.
Pemeriksaan perinciannya kita serahkan sebagai latihan.
\end{proof}

\subsection{Contoh}

\begin{example}[Aljabar]
Sebuah $k$-aljabar berdimensi hingga $A$ mempunyai perkalian $\mu: A \times
A \to A$ yang linear pada setiap entri.
(Aditivitas diperoleh dari sifat distributif gelanggang yang mendasarinya,
sedangkan perkalian skalar diperoleh dari sifat kompatibilitas).
Jadi, kita dapat memandang aljabar sebagai ruang vektor $A$
beserta pilihan tensor $t$ dalam $A^\vee \otimes A^\vee \otimes A$.

Jika $e_1,\ldots,e_n$ adalah basis bagi $A$,
maka \emph{konstanta struktur} aljabar tersebut
merupakan koleksi $\{ a^{ij}_\ell\}$ yang terdiri dari $n^3$ unsur $k$ sedemikian sehingga
\[
e_i\cdot e_j = \sum_{\ell=1}^n a^{ij}_\ell e_\ell
\]
untuk semua $1 \le i,j \le n$.
Sekarang kita dapat menuliskan tensor tersebut secara eksplisit sebagai
\[
t = \sum_{i,k,\ell} a^{ij}_\ell e_i^\vee \otimes e_j^\vee \otimes
e_\ell.
\]

Aljabar tersebut komutatif jika dan hanya jika $a^{ij}_\ell=a^{ji}_\ell$
untuk semua indeks yang mungkin. Secara ekuivalen, aljabar tersebut komutatif jika
dan hanya jika berada dalam subruang $\operatorname{Sym}^2(A^\vee) \otimes A$.
Syarat asosiativitas \emph{tidak} linear dalam konstanta
struktur. 
\end{example}

\begin{example}[Gelanggang koordinat]
Jika $V$ adalah ruang vektor, maka fungsi polinomial umum $f: V \to k$
jelas tidak linear maupun multilinear. Namun, himpunan fungsi
semacam itu secara kanonik diidentifikasi dengan $\mathcal{S}(V^\vee)$.
Hal ini sama sekali tidak mendalam, tetapi layak ditelaah sedikit.

Jika $x_1,\ldots,x_n$ adalah basis bagi $V^\vee$, maka
$f$ merupakan kombinasi linear unsur-unsur berbentuk
$x_{i_1}\cdots x_{i_d}$. Penerapan fungsi dilakukan
dengan memperluas secara linear aturan
\[
(x_{i_1}\cdots x_{i_d})(v) \mapsto x_{i_1}(v)\cdots x_{i_d}(v)
\]
untuk setiap $v \in V$.

Peta $\alpha : \mathcal{S}^d(V^\vee) \times V \to k$
dapat dipahami melalui pasangan kanonik (linear)
\[
\mathcal{S}^d(V^\vee) \times \operatorname{Sym}^d(V) \to k
\]
dan peta diagonal (nonlinear)
$\Delta: V \to \operatorname{Sym}^d(V)$
dengan $\Delta(v) = v \otimes \cdots \otimes v$.
Kemudian berlaku $\alpha(f,v) = f(\Delta(v))$.
\end{example}

\begin{example}[Bentuk kuadratik]
Bentuk-bentuk bilinear simetris pada ruang vektor berdimensi hingga $V$
dapat diidentifikasi dengan $\operatorname{Sym}^2(V^\vee)$, karena bentuk
bilinear biasa dapat diidentifikasi dengan
$(V \otimes V)^\vee=V^\vee \otimes V^\vee$.
Sebaliknya, bentuk-bentuk kuadratik pada $V$ dapat diidentifikasi dengan
$\mathcal{S}^2(V^\vee)$, karena tepat merupakan fungsi polinomial homogen berderajat $2$
pada $V$.

Berdasarkan Proposisi~\ref{prop:sym_alt_good_char},
terdapat isomorfisme kanonik di antara keduanya jika karakteristiknya
bukan $2$. Bila $b$ merupakan bentuk bilinear
simetris dan $q$ bentuk kuadratik yang bersesuaian,
peta $\operatorname{Sym}^2(V^\vee) \to \mathcal{S}^2(V^\vee)$
bersesuaian dengan
\[
q(v) = b(v,v)
\]
untuk $v\in V$, sedangkan peta \emph{polarisasi}
$\mathcal{S}^2(V^\vee) \to \operatorname{Sym}^2(V^\vee)$
bersesuaian dengan
\[
b(v,w) = \frac{1}{2}\left( q(v+w)-q(v)-q(w)\right)
\]
untuk $v,w \in V$.

Dalam kasus berdimensi dua,
bentuk bilinear simetris dengan matriks
$\begin{pmatrix} a & b\\ b&c \end{pmatrix}$ bersesuaian dengan bentuk kuadratik
$ax^2+2bxy+cz^2$. Isomorfisme dalam basis tensor bersesuaian dengan
\[
(e_1^2, e_1e_2, e_2^2) \leftrightarrow
(e_1 \otimes e_1, e_1 \otimes e_2 + e_2 \otimes e_1, e_2 \otimes e_2)
\]
untuk $e_1, e_2$ basis bagi $V$.
\end{example}

\begin{example}[Bentuk volume]
Diberikan ruang vektor berdimensi $n$, yaitu $V$,
sebuah \emph{bentuk volume} adalah unsur tak nol $\omega \in \Lambda^n(V)^\vee$.
Karena $\Lambda^n(V)$ berdimensi $1$, bentuk volume hanyalah
pilihan isomorfisme $\omega : \Lambda^n(V) \cong k$.
Diberikan basis terurut $e_1,\ldots, e_n$ dari $V$,
kita memperoleh bentuk volume yang bersesuaian
$\omega = (e_1 \wedge \cdots \wedge e_n)^\vee$.
Perhatikan bahwa urutannya diperlukan; tanpa urutan itu kita hanya mengetahui $\omega$
hingga tanda.

Bentuk volume mempunyai interpretasi geometris yang membenarkan namanya.
Dalam $\mathbb{R}^n$, terdapat bentuk volume baku
$\omega=(e_1\wedge \cdots e_n)^\vee$ dengan
$e_1,\ldots,e_n$ merupakan basis baku.
Jika $v_1,\ldots,v_n$ adalah vektor-vektor dalam $\mathbb{R}^n$,
maka $\omega(v_1\wedge \cdots \wedge v_n)$ adalah volume
paralelotop yang dibangkitkannya. 

Secara lebih umum, jika $k=\mathbb{R}$,
maka suatu bentuk volume $\omega \in \Lambda^n(V)^\vee$
bernilai positif atau negatif pada setiap basis terurut.
Karena itu, dapat dikatakan bahwa bentuk volume menentukan suatu
\emph{orientasi}.
Hal ini menggeneralisasi pembedaan antara arah jarum jam dan berlawanan arah jarum jam
untuk $\mathbb{R}^2$, serta pembedaan antara kaidah tangan kiri dan kaidah
tangan kanan untuk $\mathbb{R}^3$.
\end{example}

\begin{example}[Pasangan hasil kali baji]
Misalkan $V$ adalah ruang vektor berdimensi $n$.
Diberikan pilihan bentuk volume $\omega : \Lambda^n(V) \cong k$,
kita memperoleh pasangan sempurna
\[
\Lambda^d(V) \times \Lambda^{n-d}(V) \to
\Lambda^n(V) \cong k
\]
untuk setiap $0 \le d \le n$ dari perkalian aljabar eksterior.
Jadi, $\omega$ menentukan isomorfisme
\[
\omega_d : \Lambda^d(V) \cong \Lambda^{n-d}(V^\vee) \ .
\]
Pilihan $\omega$ yang berbeda hanya mengubah $\omega_d$ hingga suatu faktor
penskalaan, sehingga $\omega_d$ bersifat kanonik \emph{hingga penskalaan}.

Fungsi $\omega_d$ memetakan vektor-$d$ sederhana ke vektor
$(n-d)$ sederhana. Selain itu, misalkan $v_1,\ldots,v_d \in V$
dan $\ell_1,\ldots,\ell_{n-d} \in V^\vee$ memenuhi
\[
\omega_d(v_1 \wedge \ldots \wedge v_n) = \ell_1 \wedge \cdots \wedge
\ell_{n-d} .
\]
Subruang yang direntang oleh $v_1,\ldots,v_n$ tepat merupakan subruang
solusi $v$ bagi persamaan-persamaan
\[
\ell_1(v)=\cdots=\ell_{n-d}(v)=0 .
\]
\end{example}

%%%%%%%%%%%%%%%%%%
%%%%%%%%%%%%%%%%%%

\bibliographystyle{alpha}
\bibliography{rep_theory}

\end{document}


